\chapter*{Preface}
In 1827 the botanist Robert Brown watched pollen grains in water under his microscope and saw tiny particles jiggle without
rest, in a dance that never repeated and never stopped. In 1905 Einstein explained it by the collisions of molecules, and
Perrin soon used it to count them; in 1923 Wiener turned the dance into a mathematical object, a random continuous curve. This book is about that curve and its descendants: random motion in all its forms, from the path of a
single particle to the eigenvalues of a large random matrix, which move like particles that refuse to collide.

The first part builds Brownian motion and the calculus it requires. Gaussian processes and Kolmogorov's continuity criterion
give the path; L\'evy's modulus, the zero set of dimension one half and the law of the maximum describe its fine structure.
Continuous martingales and the It\^o integral make it possible to differentiate along a path that has no derivative anywhere,
and It\^o's formula, L\'evy's characterisation, Tanaka's formula and the local time are the result. Brownian motion then
solves the heat equation (Feynman--Kac), changes probability without changing the paths (Girsanov), solves the Dirichlet
problem, decides which dimensions are recurrent, and in the plane is invariant under every conformal map.

The second part puts the motion to work. Stochastic differential equations model a noisy world, from the Ornstein--Uhlenbeck
velocity to the Black--Scholes price, from Bessel processes to Dyson's Brownian motion of eigenvalues. Markov semigroups,
Poisson random measures and the L\'evy--It\^o decomposition describe every process with independent increments, jumps
included. The same part builds the branching random walk, the family tree behind the largest values of a random characteristic
polynomial. The first-order maximum is a theorem. The second-order term is stated, and marked as the frontier.

The third part is networks. A random walk on a graph and an electric current solve the same equations: recurrence is finite
resistance to infinity, and the matrix-tree theorem counts the ways to keep the network connected. Random walks on groups
lead to Kesten's theorem. Expander graphs, the semicircle and the sine kernel are spectra, and Volume~8 proves them. An appendix
keeps the comparison with the zeros of the zeta function: what the matrices suggest, and what is still a conjecture.

The aim throughout is to understand why the theorems are true, and to see them at work. Every result is proved, with one
clearly marked exception at the research frontier: the second-order term of the maximum of the characteristic polynomial.
Pause and think boxes ask you to guess
before reading on. Each chapter ends with exercises, graded from warm-ups ($\star$) to small gems ($\star\star\star$), each
with a hint and a complete solution that first explains how to find the idea and then writes it up.

The book assumes the probability of Volume~3 and Volume~10, the measure theory of Volume~6, the functional analysis of Volume~7 and
the Haar measure of Volume~8. The chapter on the largest values uses the trace moments proved in Volume~8, and the networks chapter uses the Laplacian of Appendix~A. Its beauty is in the meeting of the
continuous and the discrete: one random curve, invented to explain pollen grains, turns out to govern heat, harmonic functions,
option prices and electrical networks, and through Dyson's motion it reaches the eigenvalues whose laws Volume~8 proves, and through them,
against all expectation, the zeros of the zeta function.
